\theoremstyle{plain}
\newtheorem{proposition}{Proposition}[chapter]
\newtheorem{lemma}{Lemma}[chapter]

\theoremstyle{definition}
\newtheorem{assumption}{Assumption}[chapter]
\newtheorem{definition}{Definition}[chapter]

% Hypotheses are numbered 1, 2, 3 across the thesis and set as indented italic
% blocks, in the same style as the research question. The head stays upright bold.
\newtheoremstyle{thesishypothesis}
  {0.2\baselineskip}{0.2\baselineskip}{\itshape}{}{\upshape\bfseries}{.}{\newline}{}
\theoremstyle{thesishypothesis}
\newtheorem{hypothesisinner}{Hypothesis}
\newenvironment{hypothesis}[1][]
  {\begin{quote}\begin{hypothesisinner}[#1]}
  {\end{hypothesisinner}\end{quote}}
\crefname{hypothesisinner}{Hypothesis}{Hypotheses}
\Crefname{hypothesisinner}{Hypothesis}{Hypotheses}

\theoremstyle{remark}
\newtheorem{remark}{Remark}[chapter]
